% !TeX document-id = {84b599f9-72a4-4c88-a294-14afbd68e805}
%% Vorlage fuer Studentische Arbeiten - Institut fuer Leichtbau - RWTH
%% Erstellt von: Johanna Schlupkothen (schlupkothen@ilb.rwth-aachen.de)
%% 			und: Ulrike Schlauch (schlauch@ilb.rwth-aachen.de) 
%% Basierend auf der KOMA Script Klasse scrbook
%% Erste Version: v01-1 Stand Oktober 2013
%_______________________________________________________________________________________
% Updates:		Author 			Date			Changes
%---------------------------------------------------------------------------------------
% v01-2			SCHLUPKOTHEN	24.09.2014		- Header changed to Prof. Schroeder
% v01-3 		SCHLUPKOTHEN 	27.08.2015 		- Offizielles Institutslogo
%	 
%
%---------------------------------------------------------------------------------------
%
%---------------------------------------------------------------------------------------
%------                               --------------------------------------------------
%		!!!!!   R E A D   M E   !!!!!
%------                               --------------------------------------------------
%---------------------------------------------------------------------------------------
%
% !!!!! BEVOR IHR FRAGT, WEIL IRGENDETWAS " NICHT FUNKTIONIERT ": SCHAUT IN DIE  
% KOMMENTIERUNG DES ZUGEHOERIGEN PACKAGES IN DIESER PRAEAMBEL!!!!!!
%
% Hinweise zur Kompilierung:
% Empfohlener Editor:			TeXstudio
%
%  	Standardcompiler: 			PdfLaTeX
%  	Standardbibliographie: 		Biber
% 	Verarbeitung mit Glossar:	pdflatex % 
%								makeindex -s %.ist -t %.alg  -o %.acr  %.acn
%								makeindex -s %.ist -t %.slg1 -o %.syi1 %.syg1
%								makeindex -s %.ist -t %.slg2 -o %.syi2 %.syg2 
%								pdflatex %
% 	Definition eines eigenen Befehls fuer das Kompilieren des Glossars in TeXstudio bzw. TeXmakerX:
%	___________________________________________________________________________________________
%	->Optionen-> TeXstudio konfigurieren...->Erzeugen->Benutzerbefehle(siehe Feld unten):
%								Hinzufügen eines neuen Benutzerbefehls, bspw. user0:
%								Eintrag im 2. Feld:
%								txs:///pdflatex | txs:///makeindex -s %.ist -t %.alg -o %.acr %.acn| txs:///makeindex -s %.ist -t 	.slg1 -o %.syi1 %.syg1| txs:///makeindex -s %.ist -t %.slg2 -o %.syi2 %.syg2 | txs:///pdflatex
%	___________________________________________________________________________________________
%	->Options-> TeXstudio konfigurieren...->Befehle-> Makeindindex:
% 								Ensure that is written only: makenindex.exe
%								WITHOUT any argument like "%.idx" behind it
%
% 
%---------------------------------------------------------------------------------------
%--> Ablauf bei vollständiger Verarbeitung:
% 								pdflatex | pdflatex |biber | pdflatex | user0
%
%
%_______________________________________________________________________________________
%Steuerbefehle fuer TexShop fuer Mac (Gibt TS an, mit welchen Programmen gearbeitet werden soll)
%
% !TEX TS-program = pdflatex
% !BIB TS-program = biber
%---------------------------------------------------------------------------------------
\documentclass[	a4paper,		%
				12pt,			%
				titlepage,		%
				twoside,		%
				smallheadings,	%
				DIV				= 14,			%
				BCOR			= 6mm,			%
				parskip			= half,			%
				numbers 		= noenddot, 	%
				headings 		= small,		%
				cleardoublepage = empty]{scrbook}
%
%---------------------------------------------------------------------------------------
%			 LADEN DER P A C K A G E S
%---------------------------------------------------------------------------------------
%
\usepackage{ae}               	% almost european, virtueller T1-Font
\usepackage{array}          	% fuer aufwändigere Tabellen
\usepackage{colortbl}       	% farbige Tabellen (v. D. Carlisle)
\usepackage{longtable}      	% seitenübergreifende Tabellen
%\usepackage{ltablex}			% Kombiniert longtable mit tabularx
\usepackage{tabularx} 			% automatische Spaltenbreite
			\renewcommand{\arraystretch}{1.3}
\usepackage{color}    			% Farbiger/grauer Text
\usepackage[utf8]{inputenc}
\usepackage[british,ngerman]{babel}
\usepackage{amsfonts}
\usepackage{amsmath,amssymb}
\usepackage[sc]{mathpazo}	% Palatino Schrift
			\linespread{1.05}
\usepackage[T1]{fontenc}
\renewcommand*\familydefault{\rmdefault}
\usepackage[automark,nouppercase,headsepline]{scrlayer-scrpage}
\usepackage{pdfpages}
\usepackage[T1,hyphens]{url} 	% much like \verb allow line breaks for paths and URLs
\usepackage{calc}
%\usepackage{breqn}
\usepackage{ifthen}
\usepackage[absolute]{textpos}
\usepackage[percent]{overpic}
\usepackage{relsize}
\usepackage{lscape}
\usepackage{setspace}
\usepackage{graphicx}
\usepackage{tocbasic}
\usepackage{paralist}
\usepackage{booktabs}
\usepackage{multirow}
\usepackage{multicol}
\usepackage{rotating}
\usepackage{mathtools}
\usepackage{mdwlist}
\usepackage{anysize}
\usepackage{pslatex} 
\usepackage{wrapfig}
\usepackage{scrhack}
\usepackage{csquotes}
\usepackage{xkeyval}
\usepackage{xfor}
\usepackage{amsgen}
\usepackage{relsize}
\usepackage{xfrac}
\usepackage[cmyk,dvipsnames]{xcolor}
\allowdisplaybreaks[1]
\usepackage{tikz}				% fuer Grafiken, Baeume und Graphen
	\usetikzlibrary{decorations.pathmorphing}
	\usetikzlibrary{shapes}
	\usetikzlibrary{backgrounds} 
\usepackage{subfig}				
\usepackage[format		= hang,	%
 			margin		= 10pt,	%
 			font		= small]%
 			{caption}			% captions (Bildunterschriften) formatieren
	\captionsetup[figure]{		%
	  		name 		= Abb.} %
  	\captionsetup[table]{		%
  			name = Tabelle }	%
\usepackage{epstopdf} 			% Konvertiert eps zu pdf, wenn erkannt
\usepackage{fixmath}  			% 
\usepackage{dcolumn}  			% Ausrichtung an Komma oder Punkt
\usepackage{cancel}
\usepackage[
			   bottom,     		% Footnotes appear always on bottom. This is necessary
			               		% especially when floats are used
			   stable,      	% Make footnotes stable in section titles
			   perpage,     	% Reset on each page
			   %para,       	% Place footnotes side by side of in one paragraph.
			   %side,       	% Place footnotes in the margin
			   %ragged,      	% Use RaggedRight
			   %norule,   	  	% suppress rule above footnotes
			   multiple,  	  	% rearrange multiple footnotes intelligent in the text.
			   %symbol,     	% use symbols instead of numbers
			]{footmisc}
\usepackage{varioref} 			% Intelligente Querverweise
%
%---------------------------------------------------------------------------------------
% 			UMGEBUNG   K E Y W O R D S / S C H L A G W O R T E
%---------------------------------------------------------------------------------------
\newenvironment{keywords}%
   {\begin{trivlist}\item[]{\bfseries \iflanguage{ngerman}{Schlagwörter:}{Keywords:}}\ }% oder "Keywords:"
   {\end{trivlist}}
%
%
%---------------------------------------------------------------------------------------
%             NEW   C O M M A N D S
%---------------------------------------------------------------------------------------
%
% 
%
%---------------------------------------------------------------------------------------
% 			EINSTELLUNGEN   S E I T E N L A Y O U T
%---------------------------------------------------------------------------------------
%---------------------------------------------------------------------------------------
% 				K O P F Z E I L E   UND   F U SS Z E I L E
%---------------------------------------------------------------------------------------
% scrpage2 package
\clearscrheadfoot
\setkomafont{pageheadfoot}{\normalfont\normalcolor}
\setkomafont{pagenumber}{\normalfont}
\automark[section]{chapter}													% scrheadings mit Kapitel bzw. section und Seitenzahl oben
\ofoot[]{}																	% scrplain ohne alles
\setheadsepline{.4pt}
\ohead[]{\normalfont\pagemark}
\ihead[]{\normalfont\headmark}
% Appendix Seitenlayout
\defpagestyle{app}{{\headmark\hfill}{\headmark\hfill}{\headmark\hfill}}{ 	%
				   {\pagemark\hfill}{\hfill\pagemark}{\hfill\pagemark\hfill}}
%
%---------------------------------------------------------------------------------------
% 				STIL   Ü B E R S C H R I F T E N
%---------------------------------------------------------------------------------------
\setkomafont{sectioning}{\rmfamily\bfseries}
%
%---------------------------------------------------------------------------------------
% 				STIL   U N T E R S C H R I F T E N
%---------------------------------------------------------------------------------------
\addtokomafont{caption}{\small\normalfont}
%

%---------------------------------------------------------------------------------------
% 			ENDE   S E I T E N L A Y O U T
%---------------------------------------------------------------------------------------
%---------------------------------------------------------------------------------------
%
%---------------------------------------------------------------------------------------
% 			EINFUEGEN VON  M A T L A B   C O D E
%---------------------------------------------------------------------------------------
%
\usepackage[final]{listings}
\lstloadlanguages{Python}
\lstset{	basicstyle 			= \tiny, 									% print whole listing small
			keywordstyle 		=\color{Blue},								% für s/w:\color{black!80} keywords
			identifierstyle 	=, 											% nothing happens
			commentstyle 		= \color{Green},							% für s/w: \color{black!60} comments 
			stringstyle 		= \color{Plum},								% für s/w: black!40 typewriter strings
			showstringspaces	= false, 									% no special string spaces
			emph 				= {for, if, then, else, end},				%
			emphstyle 			= \color{Blue},								% für s/w: \color{black!80}
			%numbers 			= left, 									%  show number_line
			%numberstyle 		= \tiny, 									% style of number_line
			%firstnumber		= 1,										%
			%stepnumber 		= 5, 										% one number_line after stepnumber
			%numbersep 			= 5pt,										%
			language 			= {Python}, 								% 
			extendedchars 		= true,  									%
			breaklines 			= true,										%
			breakindent 		= 30pt,										%
			inputencoding		= latin1,									% für Umlaute im Source code
			breakautoindent		= true}
%
%
%---------------------------------------------------------------------------------------
% 			INITIALISIERUNG   B I B L A T E X
%---------------------------------------------------------------------------------------
%
\usepackage[style				= numeric,									%
			subentry,														%
			sorting 			= none,										%
			backend 			= biber]
								  {biblatex}								%
			\renewcommand{\mkbibnamelast}[1]{\textsc{#1}}
			\renewcommand{\bibname}{Bibliography}
			\DeclareFieldFormat{title}{#1\isdot}
			\DeclareFieldFormat{journaltitle}{#1\isdot}
			\DeclareFieldFormat{booktitle}{#1\isdot}
			\DeclareFieldFormat{issuetitle}{#1\isdot}
			\DeclareFieldFormat{maintitle}{#1\isdot}
			\DeclareFieldFormat{thesistitle}{#1\isdot}
			\DeclareFieldFormat[article]{title}{#1}
			\DeclareFieldFormat[inbook]{title}{#1}
			\DeclareFieldFormat[incollection]{title}{#1}
			\DeclareFieldFormat[patent]{title}{#1}
			\DeclareFieldFormat[thesis]{title}{#1}
%
%
%---------------------------------------------------------------------------------------
% 			P D F   S E T T I N G S   HYPERREF
%---------------------------------------------------------------------------------------
%
\usepackage{hyperref}
\hypersetup{pdfpagelayout 		= {TwoColumnLeft},							%
			pdftitle 			= {MA-Nummer-Jahr-ilb},						%
			pdfsubject 			= {Arbeitstitel-Kurzform},					%
			pdfcreator 			= {PDFLaTeX},								%
			pdfauthor 			= {Name Matrikelnummer},					%
			plainpages 			= false,									%
			linkcolor 			= blue,										%
			citecolor 			= purple,									%
			filecolor 			= green,									%
			urlcolor 			= blue,										%
			pdfborder 			= {0 0 0.5},								%
			raiselinks 			= true,										%
		bookmarksopenlevel		= 1,										%
			bookmarksopen		= true,										%
		bookmarksnumbered		= true,										%
			colorlinks 			= true,										%
			draft 				= false}									% draft = true fuer Druckversion
																			% draft = false fuer Displayversion (mit Hyperlinks)%
%
%---------------------------------------------------------------------------------------
% 			INITIALISIERUNG + EINSTELLUNGEN   G L O S S A R I E S
%---------------------------------------------------------------------------------------
%
\usepackage[nomain,acronym,nostyles]{glossaries}
%
% Hier werden zwei getrennte Symbolverzeichnisse angelegt: Das erste "Symbolverzeichnis"
% beinhaltet alle Lateinischen und alle Griechischen Schriftzeichen
% Das zweite Verzeichnis "Indizes" listet alle tiefergestellten Indizes auf
%
\newglossary[slg1]{notation}{syi1}{syg1}{Symbolverzeichnis} 				% Neues Vz. "Symbolverzeichnis"
\newglossary[slg2]{subscripts}{syi2}{syg2}{Indizes} 						% Neues Vz. "Indizes"
%
\newglossarystyle{symlongcolgrskip}{\renewenvironment{theglossary}{			% NEUER STIL ("symlongcolgrskip") 
	\begin{longtable}{@{}p{15mm}p{15mm}p{\dimexpr\linewidth-15mm-15mm-4	 	% FUER VERZEICHNIS "notation" 
				 \tabcolsep\relax}@{}}}{\end{longtable}}					% mit der Ueberschrift "Symbolverzeichnis"
	\renewcommand*{\glossaryheader}{ 										% basierend auf der Umgebung: longtable									
		\iflanguage{ngerman}{ 												%
			\textbf{Zeichen}&\textbf{Einheit}&\textbf{Beschreibung}}{		% Beschriftung Deutsch
			\textbf{Symbol}	&\textbf{Unit}	 &\textbf{Description}}\\		% Beschriftung Englisch
		\midrule\endhead}													% 
	\renewcommand*{\glsgroupheading}[1]{}									% keine Beschriftung zw. Gruppen 
	\renewcommand*{\glossaryentryfield}[4]{									% Haupteintraege in einer Zeile:  
		\glstarget{##1}{##2}&{##4}		 & ##3\\}							% Symbol / Einheit / Beschreibung
	\renewcommand*{\glsgroupskip}{	&		 &\\}							% nichts zw. Gruppen 
}
%
\newglossarystyle{symlongcol}{\renewenvironment{theglossary}{ 				% NEUER STIL ("symlongcol") 
	\begin{longtable}{@{}p{15mm}p{15mm}p{\dimexpr\linewidth-15mm-15mm-4 	% FUER VERZEICHNIS "subscripts" 
				 	 \tabcolsep\relax}@{}}}{\end{longtable}}				% mit der Ueberschrift "Indizes"
	\renewcommand*{\glossaryheader}{										% basierend auf der Umgebung: longtable
		\iflanguage{ngerman}{ 												%
			\textbf{Zeichen}&	 			 &\textbf{Beschreibung}}{		% Beschriftung Deutsch
			\textbf{Symbol}	&				 &\textbf{Description}}\\		% Beschriftung Englisch
		\midrule\endhead}													% 
	\renewcommand*{\glsgroupheading}[1]{}									% keine Ueberschriften zwischen Gruppen  
	\renewcommand*{\glossaryentryfield}[4]{									% Haupteintraege in einer Zeile:  
		\glstarget{##1}{##2}& {##4}		& ##3\\}							% Symbol / Einheit / Beschreibung 
	\renewcommand*{\glsgroupskip}{}											% nichts zw. Gruppen 
}
% 
\newglossarystyle{acrolongcol}{\renewenvironment{theglossary}{				% NEUER STIL ("acrolongcol") für
	\begin{longtable}{@{}p{30mm}p{\dimexpr\linewidth-30mm-4 				% \acronymtype mit der 
					 \tabcolsep\relax}@{}}}{\end{longtable}}				% Ueberschrift "Abkürzungsverzeichnis"
	\renewcommand*{\glossaryheader}{}										% Ohne Beschriftung
	\renewcommand*{\glsgroupheading}[1]{}									% Ohne Gruppenueberschrift
	\renewcommand*{\glossaryentryfield}[4]{									% Eintraege:
		\glsentryitem{##1}\glstarget{##1}{##2}	&##3.\\}					% Kurzform / Langform
																			% mit numberedlist: ##3. \hspace{2ex}(##5)
	\renewcommand*{\glsgroupskip}{} 
}																% Glossar-Befehle anschalten
%
%---------------------------------------------------------------------------------------
% 			INITIALISIERUNG   M A K E I N D E X
%---------------------------------------------------------------------------------------
%
\usepackage{makeidx}														% Index erstellen
\usepackage{float}															% neue float-Umgebungen definieren, aus 
																			% Kompatibilitaetsgruenden ganz zum Schluss einbinden
\makeglossaries
%
%---------------------------------------------------------------------------------------
% 			LADEN   E I N T R A E G E   ABKUERZUNGSVERZEICHNIS + LITERATURVERZEICHNIS
%---------------------------------------------------------------------------------------
\loadglsentries{Input_glossaries}											% .tex Dateiname mit Nomenklatur-Eintraegen
\bibliography{Biblio}														% *.bib Dateiname mit Literatureintraegen
%
%
%---------------------------------------------------------------------------------------
% 			EINSTELLUNG   O P T I O N E N
%---------------------------------------------------------------------------------------
%\setcounter{tocdepth}{0}													% nur bis 1.1/1.2/... ins Inhaltsvz.
%\setcounter{secnumdepth}{0}   												% Kapitel-Nummerierung nur bis 1.1.1/ 
\setcounter{secnumdepth}{5}
\setcounter{tocdepth}{5}

\clubpenalty 					= 10000										% Schusterjungen und Hurenkinder verhindern.
\widowpenalty 					= 10000										% Wen's interessiert: 
\displaywidowpenalty 			= 10000										% http://de.wikipedia.org/wiki/Hurenkind_und_Schusterjunge
\interlinepenalty 				= 5000
\relpenalty						= 5000
%
%
%---------------------------------------------------------------------------------------
%========  B E G I N N = D O C U M E N T  =============================================%
%---------------------------------------------------------------------------------------
\begin{document}
\deffootnote[1.5em]{1.5em}{1em}{\textsuperscript{\thefootnotemark}}
%
% 			TITLEPAGE
%---------------------------------------------------------------------------------------
%\includepdf{titelseite.pdf}	\cleardoublepage 							% Einbinden der Deckseite (pdf von Matsen erstellt)
\frontmatter
\pagestyle{empty}															% keine Kopf/Fusszeile, keine Seitenzahlen
	%s%%% Titelseite
\begin{titlepage}
\begin{textblock*}{110mm}(88mm,12mm)
\includegraphics*[height=16mm]{./pictures/rwth_sla}
\end{textblock*}
\begin{textblock*}{\textwidth}(30mm,60mm)
\centering
{\large\bfseries Projektarbeit} \\[3ex]
\begin{huge}\bfseries\scshape Studie zur Untersuchung des Verhaltens von Low-Velocity Impacts unter verschiedenen Parametern\par\end{huge}
\end{textblock*}
\vspace*{85mm}
\centering\large
Diese Arbeit wurde vorgelegt am \\[3ex]
Institut für Strukturmechanik und Leichtbau (SLA)\\
Univ.-Prof.~Dr.-Ing. Kai-Uwe Schr\"{o}der\\
Fakultät für Maschinenwesen\\
RWTH Aachen\\[5ex]
von\\[1ex]
\textbf{Julian Callard}\\
\textbf{Finn Eggers}\\
\vskip 5ex
Betreuer: \\[1ex]
%\begin{tabular}{p{.42\textwidth}p{.05\textwidth}p{.42\textwidth}}
%Dipl.-Ing. Max Musterbetreuer						& $\:$	&	Dr.-Ing. Emil Externbetreuer\\
%{\small Institut für Strukturmechanik und Leichtbau}&		& 	{\small Abteilung für Schnickschnack}\\[-.5ex]
%{\small RWTH Aachen	University}						& 		&	{\small Musterunternehmen, Musterstadt}
%\end{tabular}
Stephan Kalapis, M.Sc.\\
SLA, RWTH Aachen\\[1ex]

\begin{textblock*}{\textwidth}(30mm,267mm)
\centering\large Aachen, Juli 2020
\end{textblock*}

\end{titlepage}
%% Titelseite Ende

			\cleardoublepage 					
%
% 			SUBJECT + EIGENSTAENDIGKEIT
%---------------------------------------------------------------------------------------
\pagestyle{empty}															% keine Kopf/Fusszeile, keine Seitenzahlen
 	
 	\selectlanguage{ngerman}												% Erzwungen: Deutsch fuer Eigenstaendigkeitserklaerung
 	
 	{\Huge \bf Sperrvermerk} \\

\vspace*{1cm}

\begin{figure}
	\includegraphics[width=\linewidth]{Pictures/Pictures/PorscheSil.png}
\end{figure}	

Der vorliegende Praktikumsbericht beinhaltet interne und teilweise vertrauliche Daten
des Unternehmens: Dr. Ing. h.c. F. Porsche AG.\\

Die Einsicht in den Praktikumsbericht ist Unbefugten nicht gestattet. Ausgenommen
hiervon sind die Gutachter, sowie berechtigte Mitglieder des Prüfungsausschusses.\\
Die Vervielfältigung und Veröffentlichung des Praktikumsberichts – auch auszugsweise – ist
grundsätzlich nicht erlaubt. Eine Ausnahme von dieser Regelung bedarf einer Erlaubnis
der Dr. Ing. h.c. F. Porsche AG.\\



%I hereby assure that I composed the thesis at hand self-reliantly. 







	\cleardoublepage
 	%{\footnotesize \begin{center}-- Diese Titelseite ist nicht die Originaltitelseite zur Abgabe der Arbeit --\end{center}}
%---------------------------------------------------------------------------------
% Institutskopf
%---------------------------------------------------------------------------------
%\begin{tabularx}{\textwidth}{lXr}
%\parbox{100mm}{\centering\bf LEHRSTUHL UND INSTITUT F{\"U}R LEICHTBAU\\ %
%	 			Rheinisch-Westf{\"a}lische Technische Hochschule Aachen\\ %
%	 			Universit{\"a}tsprofessor Dr.-Ing. K.-U. Schr\"{o}der}	&&%
%\parbox{35mm}{\small D-52062 Aachen\\ W{\"u}llnerstr. 7\\ Tel: 0241/80-96830\\ Fax:  0241/80-92230}
%\end{tabularx}

%
%---------------------------------------------------------------------------------
% Titel
%---------------------------------------------------------------------------------
\newpage
\begin{textblock*}{\textwidth}(30mm,60mm)
\centering
{\bfseries Praktikumsbericht} \\[1ex]
\bfseries\scshape Erklärung zur selbstständigen Anfertigung des Berichts\par
\end{textblock*}
\vspace{5ex}

\begin{textblock*}{\textwidth}(30mm,100mm)
\centering	
Ich versichere durch meine Unterschrift, dass ich diese Arbeit selbstständig angefertigt
und keine anderen als die angegebenen Hilfsmittel benutzt habe. Wörtlich oder
sinngemäß aus anderen Quellen übernommene Textstellen, Bilder, Tabellen u. a. sind
unter Angabe der Herkunft kenntlich gemacht.\par
\end{textblock*}

\vspace{5ex}

\selectlanguage{ngerman}
\vspace*{10cm}\begin{tabular}{p{7cm}p{5cm}}
	Stuttgart, den \today & $\rule{5cm}{0.01cm}$ \\
	& (Unterschrift, Julian Callard)
\end{tabular}


\newpage

			\cleardoublepage							% genstaendigkeitserklaerung
%
% 			PREFACE + ABSTRACT
%---------------------------------------------------------------------------------------
\pagestyle{scrheadings}														% Koma-script Kopf-/Fusszeile , Definition s.o.
%\pagenumbering{roman}														% Seitennummerierung in kleinen roemischen Zahlen
	%\chapter*{Kurzfassung} 
\markboth{Kurzfassung}{Kurzfassung}
%% ==============================


Die Schadensbeurteilung von Platten, auf die ein Objekt gefallen ist oder ein Stoß verübt wurde, stellt eine wichtige Aufgabe in der Industrie dar. Jedoch werden in den meisten Arbeiten lediglich die Auswirkungen eines Stoßes auf eine Platte in der Plattenmitte analysiert. Die Position des Stoßes auf die Platte hängt allerdings eng mit der Schadensbeurteilung zusammen. Im Folgenden wird eine Parameterstudie durchgeführt, welche Einwirkungen wie die Position, Gewicht, Geschwindigkeit und die Platteneigenschaften kombiniert. Zunächst wird festgestellt, unter welchen Bedingungen die Berechnungen gültig und zuverlässig sind. Im Anschluss wird die eigentliche Parameterstudie anhand eines Programms durchgeführt, welches die zuvor dargestellten mathematischen Grundlagen verwendet.
\begin{keywords}
Stoßvorgang, Biegung, Numerik, isotrope Materialien, Platten
\end{keywords} 	\cleardoublepage
	\selectlanguage{british}												% Erzwungen: Englisch fuer Abstract
	%\chapter*{Abstract} 
\markboth{Abstract}{Abstract}
%% ==============================

The damage assessment of plates on which an object has fallen or an impact has been made is an important task in the industry. However, in most work, only the effects of an impact on a plate in the middle of the plate are analysed. However, the position of the impact on the plate is closely related to the damage assessment. In the following, a parameter study is carried out which combines effects such as position, weight, speed and panel properties. First, it is determined under which conditions the calculations are valid and reliable. Then the actual parameter study is carried out using a program that uses the mathematical principles described above.

\begin{keywords}
\LaTeX, impact, bending, numerics, isotropic materials, plates
\end{keywords}			\cleardoublepage
%	\selectlanguage{ngerman}												% Erzwungen: Deutsch fuer Vorwort
%	\include{preface}			\cleardoublepage							% Vorwort
%
% 			LISTINGS
%---------------------------------------------------------------------------------------
	\tableofcontents			\cleardoublepage							% Inhaltsverzeichnis
%	\listoffigures 				\cleardoublepage							% Wollte der Chef streichen
%	\listoftables				\cleardoublepage							% Wollte der Chef streichen
	%\printglossary[	type 		= \acronymtype,								% 
	%				style 		= acrolongcol,								% 
	%				title		= Abkürzungsverzeichnis,					% 
	%				nonumberlist] \cleardoublepage							%  
																			% 
	%\printglossary[	type 		= notation,									%
	%				style 		= symlongcolgrskip,			 				%
	%				nonumberlist] \clearpage
	%\printglossary[	type 		= subscripts,								%
	%				style 		= symlongcol, 								%
	%				nonumberlist] \cleardoublepage

%
% 			CONTENT
%---------------------------------------------------------------------------------------
\mainmatter
\pagestyle{scrheadings}														% Koma-script Kopf-/Fusszeile , Definition s.o.
%\pagenumbering{arabic}														% Seitennummerierung in arabische Zahlen
	\chapter{Die Porsche AG}
\label{chap:Intro}

Porsche wurde 1931 von Ferdinand Porsche gegründet. Seit 1945 stellt die Firma Sportwagen und Luxusautos her.\\
Im Finanzjahr 2018 erzielte der Konzern einen Umsatz von 25,784 Mrd. Euro.\\ %\cite{hier quelle wikipedia einfügen}
Die aktuelle Produktpalette der AG umfasst für Serienfahrzeuge:\\

\begin{itemize}
	\item{Macan (2013)}
	\item{Cayenne (2002)}
	\item{911 (1963)}
	\item{Boxter (1996)}
	\item{Panamera (2009)}
 \end{itemize}

\section*{Weissach}

Seit 1971 werden diese Fahrzeuge am Standort Weissach entwickelt. Das sogenannte "Entwicklungszentrum Weissach", kurz EZW, auf dem die Porsche AG ein Fahrzeug entwickeln, testen und justieren kann, bevor es in Serienproduktion in Zuffenhausen oder einem anderen Standort geht. \\
%Hier Bild vom EZW einfügen
Alleine an diesem Standort beschäftigt Porsche um die 6500 Mitarbeiter und Mitarbeiterinnen. Zuletzt wurde hier auch der Porsche 918 entickelt: Ein Supersportwagen mit hybridem Antrieb und einem Listenpreis von mehr als 800 000 Euro.\\
%hier halb-seitiges Bild vom 918
Das in Weissach entwickelt wird, sieht man schon an der Anzahl an Patenten die hier jährlich angemeldet werden. Knapp 400 pro Jahr sind es aktuell, mit mehr als 7000 bereits zugelassenen Patenten.\\
Die Signifikanz des Standortes Weissach wird immer wieder durch spezielle Produktvarianten oder Ausstattungen hervorgehoben. Ein Beispiel hierfür ist das Weissach-Paket, welches man für den 911 bestellen kann. Teil hiervon sind typischerweise die Farben Grün und Weiß, zusammen mit besonderen Aero-Elementen.

%\section*{Porsche als Ausbildungsbetrieb}


		\cleardoublepage
	\chapter{Das Praktikum}

\section{Fachpraktikum}

Für das Fachpraktikum mussten noch 14 Wochen absolviert werden. Alle Bereiche des Grundpraktikums wurden bereits im Vorpraktikum abgearbeitet, weshalb sich das Fachpraktikum nur auf das in den Praktikumsrichtlinien der RWTH Aachen genannte Fachpraktikum Teil A und Teil B bezieht.\\
%\cite{Praktikumsrichtlinien der RWTH Maschbau}
Die 8 Wochen aus Teil B konnten alle mit FP7 abgegolten werden, da in Weissach ausschließlich Fahrzeugentwicklung betrieben wird. Die übrigen sechs Wochen aus Teil A wurden mit FP4 - \textbf{Messen, Prüfen, Qualitätskontrolle} - mit einem Anteil von drei (3) Wochen und FP6 - \textbf{Montage} - mit weiteren drei Wochen ausgeführt.\\

\section{Fachbereich Porsche}

Bei Porsche wurde das Praktikum im Fachbereich \textbf{Aerodynamik und Thermomanagement} absolviert. \\
Diese Abteilung ist für die thermische Absicherung der Fahrzeuge und die aerodynamische Auslegung der Neuentwicklungen zuständig.\\
Neben Versuchen im Klimawind- oder Windkanal werden auch Erprobungsfahrten, Simulationen und andere Mess- und Testmethoden verwendet um die hohen Standards bei Porsche aufrecht zu erhalten. \\
Im engen Austausch mit den anderen Abteilungen und Bauteilverantwortlichen wird so jedes Fahrzeug und Einbauteil erprobt und entsprechend der geltenden Normen und Rechtssprechungen, sowie der Anforderungen die Porsche an sich selbst stellt, entwickelt und zur Serienreife gebracht.

		\cleardoublepage
	%\chapter[Implementierung (eigene Leistung)]{Implementierung}
\label{chap:Implementation}


Das Python-Skript folgt dem dargestellten Schema in \ref{fig:flowchart}. Zunächst werden Parameter der Funktion \texttt{compute()} übergeben. Wenn keine Parameter übergeben werden, wird direkt der Basisfall berechnet. Da die Größe $S(k)$ die meiste Zeit zum Berechnen benötigt, sie jedoch nur von der Geometrie der Platte abhängt, kann für viele aufeinanderfolgenden Rechnungen das $S(k)$ aus vorherigen Rechnungen übernommen werden.\\ 
Im Anschluss wird der erste Zeitschritt separat von den darauffolgenden Zeitschritten berechnet. Da eine numerische Integration durchgeführt wird, wird zunächst eine Annäherung berechnet, welche ebenfalls für eine zweite Näherung genutzt wird.\\
Es ist theoretisch möglich weitere Annäherungen durchzuführen. Es hat sich jedoch gezeigt, dass die Fehler bereits nach der zweiten Annäherung verschwinden.


\begin{figure}[!h]
	\begin{center}
		\begin{overpic}[width=\linewidth]{pictures/FlowChart.eps}
			\put(29,59){\scriptsize{Start}}
			
			\put(18,35){\scriptsize{Nein}}
			
			\put(33,25){\scriptsize{Ja}}
			
			\put(70,25){\scriptsize{Nein}}
			
			\put(70,45){\scriptsize{Ja}}
			
			
			\put(-19,46){\begin{minipage}{\textwidth}\centering{\scriptsize{Konstanten}} \\ \scriptsize{bestimmen}\end{minipage}}
			
			
			\put(-19,16.6){\begin{minipage}{\textwidth}\centering{\scriptsize{1. Zeitschritt:}} \\ \scriptsize{1. Näherung für} \\ \scriptsize{$u,v,w,P,z$} \end{minipage}}
			\put(-19,5.2){\begin{minipage}{\textwidth}\centering{\scriptsize{1. Zeitschritt:}} \\ \scriptsize{2. Näherung für} \\ \scriptsize{$u,v,w,P,z$} \end{minipage}}
			%\put(23.5,16.6){\begin{minipage}{\textwidth}1. Näherung \\ für $u,v,w,P,z$\end{minipage}}
			%\put(23.5,5){\begin{minipage}{\textwidth}2. Näherung \\ für $u,v,w,P,z$\end{minipage}}
			
			%\put(-19,34){\begin{minipage}{\textwidth}\centering{$S(k)$} \\ gegeben?\end{minipage}}
			\put(17.5,33.5){\begin{minipage}{\textwidth}\centering{\scriptsize{Zeitschritte}} \\ \scriptsize{übrig?}\end{minipage}}
			\put(-19,33.5){\begin{minipage}{\textwidth}\centering{\scriptsize{$S(k)$}} \\ \scriptsize{gegeben?}\end{minipage}}
			
			
			\put(-42,25){\begin{minipage}{\textwidth}\centering{\scriptsize{$S(k)$}} \\ \scriptsize{bestimmen}\end{minipage}}
			
			
			\put(42,57.5){\begin{minipage}{\textwidth}\centering{\scriptsize{n. Zeitschritt:}} \\ \scriptsize{1. Näherung für} \\ \scriptsize{$u,v,w,P,z$} \end{minipage}}
			\put(42,46){\begin{minipage}{\textwidth}\centering{\scriptsize{n. Zeitschritt:}} \\ \scriptsize{2. Näherung für} \\ \scriptsize{$u,v,w,P,z$} \end{minipage}}
			
			%\put(84.5,57.5){\begin{minipage}{\textwidth}\small{1. Näherung} \\ \small{für $u,v,w,P,z$}\end{minipage}}
			%\put(84.5,46){\begin{minipage}{\textwidth}2. Näherung \\ für $u,v,w,P,z$\end{minipage}}
			
			\put(17,5){\begin{minipage}{\textwidth}\centering{\scriptsize{Ergebnisse}} \\ \scriptsize{auswerten}\end{minipage}}
			
		\end{overpic}
		\caption{Flowchart des Python-Skriptes}
		\label{fig:flowchart}
	\end{center}
\end{figure}


\newpage

Anschließend wird eine äquivalente Rechnung für alle weiteren Zeitschritte durchgeführt und im Anschluss gegebenenfalls eine graphische Darstellung der Durchbiegung und Kraft angezeigt sowie die maximale Kraft und Durchbiegung bestimmt.\\
Die Funktion  \texttt{compute()} gibt ein Touple \texttt{time, j, tau, w, P, u, cos} zurück.  \texttt{time} ist eine Liste mit den zugehörigen Zeitpunkten zu den Durchbiegungen \texttt{w} und den Kräften \texttt{P}. \texttt{u} stellt zusätzlich die Position des Impaktors dar. Da es möglich ist, die Berechnung durch \texttt{stop\_computation\_after\_first\_impact = True} vorzeitig zu beenden, nachdem der erste Schlag beendet wurde, ist die Anzahl der Zeitschritte als \texttt{j} wiedergegeben. Die verwendete Schrittweite wird als \texttt{tau} zurückgegeben.\\
Für eine Wiederverwendung von $S(k)$ wird dieses als Liste zurückgegeben, mit dem Namen \texttt{cos}.
Der verwendete Code ist im Anhang~\ref{chap:Anhang_Skript} zu finden.

Für jeden Zeitschritt werden 2 Näherungen durchgeführt. Theoretisch können auch mehr als zwei Iteration pro Zeitschritt durchgeführt werden, jedoch hat dies keine wesentliche Änderung im Ergebnis gezeigt.

Im ersten Zeitschritt wird angenommen, dass die Kraft in jenem Zeitschritt ungefähr der Kraft im vorherigen Zeitschritt entspricht. Hieraus lässt sich direkt die neue Geschwindigkeit und Durchbiegung der Platte annähern. Ebenfalls folgt aus den zuvor genannten Größen die Durchdringung in die Platte, welche genutzt wird, um im zweiten Zeitschritt eine bessere Näherung für die Kraft zu finden. Hierfür wird das Stoßgesetz verwendet.

Um die Rechnungen auszuwerten, wird für jede Variable eine Schleifenstruktur implementiert. Der $S(k)$-Term muss nur neu berechnet werden, wenn sich die Geometrie der Platte ändert. Also wird mit zwei \texttt{for}-Schleifen iteriert. Die erste Schleife iteriert über den zu ändernden Parameter und errechnet $S(k)$, dann iteriert die zweite Schleife innerhalb der ersten mit konstantem Parameter über die Menge der $MR$.  Daraus ergibt sich eine wesentlich kürzere Rechenzeit, da das Erstellen von $S(k)$ mit Abstand am meisten Zeit in Anspruch nimmt. \\
Die Ergebnisse werden letztendlich in eine \texttt{.dat}-Datei geschrieben. 


\section*{Numerische Instabilität}

Das Skript weist gewisse numerische Instabilitäten auf. Einerseits führt eine zu geringe Plattendicke zu einer geringen Steifigkeit. Dadurch ist bereits im ersten Zeitschritt die Durchbiegung der Platte größer als die zurückgelegte Strecke des stoßenden Körpers. Dies führt zu einer negativen Durchdringung und ist somit nicht plausibel. Dieses Problem kann jedoch durch eine deutlich feinere Schrittweite behoben werden.






	\cleardoublepage
	\chapter{Praktikumsbericht}
\label{Bericht}

\section{FP7 - Entwicklung (8 Wochen)}
\label{FP7}

Der Anfang des Praktikums wurde mit Einarbeiten verbracht.\\
Zu den verwendeten Tools gehören unter anderem PowerFlow, PowerTherm, ANSA, Star CCM+ und MATLAB.\\

%bild von powerflow einfügen

Vor allem PowerFlow spielt eine große Rolle, da es für Fahrzeugsimulationen geeignet ist. Das Programm baut auf der \textbf{Lattice-Boltzmann-Methode} auf.\\

\subsection{Die Lattice-Boltzmann-Methode (LBM)}
\label{LBM}
Die Lattice-Boltzmann-Methode beruht auf einer Vereinfachung der Euler-Gleichung und ist Aufgrund geringer Speicheranforderungen pro Zelle geeigenet um komplexe Geometrien zu berechnen. \\
Da Fahrzeuge nur selten bis in den kompressiblen Machzahlbereich von 

\begin{equation}
		\label{MachEq}
	Ma > 0.3
\end{equation}

ausgelegt werden, können die getroffenen Vereinfachungen angewendet und dadurch die \textbf{LBM} für die Fahrzeugauslegung verwendet werden. \\
Da die LB-Methode leicht für die Wärmeübertragung erweitert werden kann, spielt bei der Auslegung der Fahrzeuge auch PowerTHERM eine große Rolle für die Bauteilsicherung.

\subsection{Einarbeitung}
\label{Einarbeitung}

Für die Einarbeitung in PowerTherm und PowerFlow wurden Tutorien durchgearbeitet.\\
Um sich mit den Funktionalitäten des Programms vertraut zu machen, wurde zuerst eine einfache NACA-Flügelgeometrie geladen und diese dann mit den entsprechenden Netzen und Berechnungsvolumen ausgestattet um eine Berechnung zu ermöglichen. \\
Anschließend mussten die Randbediungen an den Simulationsvolumengrenzen definiert werden und die Strömungskonditionen definiert werden. Hierzu zählen die Reynoldszahl, Temperatur, Strömungsgeschwindigkeit und Vieles mehr.\\

Nachdem die Berechnung abgeschlossen ist, wird dann PowerCASE, ein Programm zur Darstellung der Simulationsergebnisse, verwendet um sich verschiedene relevante Parameter oder Geschwindigkeiten, Drücke und vieles mehr anzeigen zu lassen. 

%hier overview von powerflow UI
Nach dem NACA-Profil wurde als nächstes mit ANSA, einem Pre-Processing Tool, gearbeitet. In diesem können Geometrien für die Simulationen aufarbeitet und vernetzt werden. \\
%hier ANSA einfügen
Abgesehen von der recht veralteten Benutzeroberfläche ist ANSA was die Geometrievorbereitung für PowerFlow gut geeignet, da die Geometriegüte in PowerFlow für verwertbare Ergebnisse nicht perfekt sein muss. \\
In ANSA wurde nun eine einfache Geländewagengeometrie aufgebaut, die dann in PowerFlow geladen und mit einem einfachen Simulationsvolumen berechnet wurde. Anschließend wurden die Ergebnisse wieder in PowerCASE ausgewertet und mit anderen Praktikanten verglichen.\\

Wenn man PowerFlow mit Star CCM+ vergleicht kommt heraus, dass CCM+ zwar etwas genauere Ergebnisse liefert. Dafür benötigt es jedoch eine sehr viel höhere Güte an eingeladener Geometrie, was es für die aktuelle Anwendung zum Teil ungeeignet macht. \\
%pros and cons of star and pwoerflow
%make a comparison vector graph in Stuggi

Am Ende der Einarbeitung wurde noch eine Einweisung in die Fahrzeughandhabung auf dem Werksgelände durchgeführt. Da die Fahrzeuge selbst auf der niedrigsten Austattungsstufe mehr \textbf{kW} entwickeln als die meisten andere Kraftwagen, ist hier besondere Vorsicht geboten. Vor allem, wenn man Fahranfänger ist. \\

\newpage
\subsection{Datenauswertung in Excel}
\label{Excel}

Als erstes Projekt wurde eine Datenauswertung in Excel zugewiesen. \\
Um mit den großen Datenmengen umzugehen wurde in Excel mit Hilfe des integrierten VBA-Editors ein Makro geschriebe. Dieser kann nach minimaler Vorarbeit des Anwenders automatisch die Datenmengen auswerten und das entsprechende Kollektiv erstellen. \\
Die hinterlegte Logik und Datenstrukturen mussten zuerst entworfen und dann implementiert und getestet werden. \\
Da VBA (Virtual Basic) eine free-type Programmiersprache ist, können Datentypen direkt im Code gecasted und definiert werden. Dies ist im Vergleich zu einer Sprache wie JAVA zwar schneller zum Tippen, ist aber bei der Code-Compilation und späteren Ausführung langsamer, da der Interpreter die Datentypen bei der Ausführung zuweisen muss. \\
Ein weiteres Problem, das speziell bei der Anwendung von VBA in der Microsoft Office-Umgebung entsteht, ist die extrem langsame Zugriffszeit von dem VBA Editor in das entsprechende Dokument in z.B. Excel.\\

\subsubsection{Makroversion 1.0}
\label{Makro1.0}

In der ersten Version des Makros wurde in den entsprechenden Iterationsschleifen direkt auf das jeweilige Worksheet in Excel zugegriffen um die Messdaten auszulesen und vergleichen zu können. \\
Die grundliegende Logik ist in Abbildung ()()() zu sehen.
%\ref{}
Da die Messdatenmengen des Versuchs aus mehreren tausend Messpunkten besteht, brauchte der Makro für die Iteration über alle Datenpunkte geschätzte 120 Stunden. Diese relativ lange Bearbeitungszeit war zwar zum Teil auf die Hardware zurückzuführen, wurde jedoch vor allem von dem oben beschriebenen Problem hervorgerufen. \\
Bei einer Messungsreihe mit 5 verschiedenen Messumgebungen die aus je 12 Messtemperaturen bestehen wird die Datenmenge schnell sehr groß. Die Messdaten wurden auf eine uniforme Zeitbasis von 1 Hz bestimmt. Je nach Messumgebung variierte auch die Messdauer. In der Testumgebung wurden daher nur etwa 20 000 Datenpunkte ausgewertet.\\

Da selbst in der Testumgebung die Berechnungsdauer bei mehreren Stunden lag, musste die Berechnungszeit verkürzt werden.

\subsubsection{Makroversion 1.5}
\label{Makro1.5}

Nach erweiterte Recherche konnte die Berechnungszeit verkürzt werden, indem die Datenauslegung am Anfang der Iteration einmal durchgeführt wurde und dann nur noch mit den eingelesenen Daten direkt in der VBA Umgebung gerechnet wurde. \\
Hier war nun vor allem die Indizierung der Matrizen eine Herausforderung. Je nachdem ob bestimmte Messinformationen mit ausgewertet werden mussten, mussten die Indizes und die Iterationsschrittweite angepasst werden. \\
Dies wurde durch eine VBA Userform gelöst, in welcher der Benutzer angeben konnte ob die Information benötigt wird. Im Code wurde basierend auf dieser Auswahl eine Iterationsvariable und die Iterationsgrenzen angepasst. \\
Nach weiterem Testen war diese Art der Auswertung erfolgreich. \\

Nun musste nur noch das Problem der Matrixindizierung gelöst werden. Um den Makro möglichst benutzerfreundlich zu gestalten, mussten die Matrixgrenzen so einprogrammiert werden, dass die Auswertung auch funktioniert wenn der Benutzer die Daten falsch in Excel einfügt oder Daten fehlen. Hierfür mussten in die Iterationsschleifen Logikabfragen implementiert werden, da die Fehlerbehandlung in VBA fragwürdig ist. Diese Maßnahme verlangsamten die Auswertung wieder etwas. Jedoch war es möglich durch Deaktivieren der Bildschirmanpassung und Hintergrundberechnung von Excel selbst diese Rechenzeit wieder auszugleichen. \\

\subsubsection{Makroversion 2.0}
\label{Makro2.0}

Nachdem die Matrixauswertung erfolgreich war, mussten die Auswertungsergebnisse noch aufgearbeitet werden. Die Ergebnismatrix wurde dann in ein eigenes Worksheet geschrieben. In diesem Worksheet ist Code hinterlegt, der die Matrixdaten in einer 3D-Grafik wiedergibt. \\
Mit dieser Matrix und dem entsprechenden Graphen können dann die Gespräche mit den Zulieferern und Lebenszeitexperten geführt werden im das Produkt möglichst optimiert zu entwickeln.\\

Durch die oben genannten Optimierungen konnte die Berechnungsdauer in der finalen Version des Makros auf etwa 20 Minuten reduziert werden. Diese Zeit wurde in der echten Auswertungsumgebung mit ca. 70 000 Datenpunkten erreicht.\\

Ein weiterer Grund warum diese Auswertung jedoch in Excel nur bedingt Effizient ist, ist der pre-processing Aufwand um verwertbare Ergbenisse zu erlangen. Die Daten müssen zuerst mit anderen Programmen aus der Messdatei ausgelesen und dann in Excel exportiert werden. Aus dieser Excel-Datei müssen die relevanten Daten dann in das Excel-Workbook mit dem hinterlegten Makro übertragen werden. Bei insgesamt 60 Messläufen ist dieser manuelle Aufwand wesentlich größer als die später benötigte Zeit um die Auswertung durchzuführen.\\

\subsubsection{Das Excel-Workbook}
\label{Workbook}

Das Excel-Workbook in welchem der Makro hinterlegt ist wurde als Vorlage (Template) erstellt und in einem entsprechendem Ordner abgespeichert um den restlichen Abteilungsmitgliedern den Zugriff auf das Workbook zu ermöglichen.\\

Da in der Fahrzeugherstellung viele verschiedene Umweltzustände und dadurch Temperaturen dargestellt und getestet werden müssen, sind in dem Excel-Workbook verschiedene Worksheets hinterlegt, in denen die relevanten Umgebungsdaten hinterlegt sind. \\
Um diese in der Berechnung automatisch mit zu Berücksichtigen müssen diese vor Starten des Makros entsprechend über eine User-Form ausgewählt werden. Um die Bearbeitung zu erleichtern ist als Standard-Wert der am häufigsten verwendete Fall hinterlegt. \\

Die Lebensdauerfaktoren werden in dem ersten Worksheet festgelegt und berechnet. Der Makro liest sich diese Werte entsprechend aus und verwendet sie. Daher ist es wichtig, dass das Makro-Template immer den neuesten Werten entsprechend angepasst wird. Diese Arbeit muss leider weiterhin manuell ausgeführt werden, da die Lebensdauerfaktorenberechnung oft in Arbeitskreisen durchgeführt wird und diese Faktoren dann nicht in einer zentralen Datenbank hinterlegt werden. \\

\newpage
\subsubsection*{(Notiz)}
\label{Note1}

In diesem Bereich ist nach Meinung des Authors weitere Entwicklung gewünscht um Arbeitsprozesse zu Streamlinen und die gesamte Entwicklung effizienter zu gestalten.

\subsubsection*{(Notiz-Ende)}

Für die Weiterverarbeitung der berechneten Daten ist es nun günstiger, die resultierenden Matrizen aus dem Makro-Worksheet in ein Worksheet ohne hinterlegte Makros und Berechnungsparametern zu transferieren. Dies resultiert aus der Geheimhaltungsklausel der Porsche AG. Die Makro-Datei hat unter anderem konzern-interne Daten hinterlegt, die nicht von Zulieferern oder Dritten Parteien eingesehen werden dürfen.\\
Hier könnte der Makro erweitert werden, damit die Ergebnisse automatisch in ein neues Excel-Workbook übertragen werden.\\

Ein weiterer Mangel an dem Marko-Workbook besteht darin, dass während der Makro läuft Excel unbenutzbar auf diesem Rechner wird. Der von Excel blockierte Thread wird komplett von dem Makro übernommen und lässt das Betriebssystem Excel als ein Programm ohne aktive Rückmeldung anzeigen. \\
Dies muss leider mit der aktuellen Version des Makros hingenommen werden.\\
Aufgrund des Wunsches diesen Workflow in MATLAB effizienter zu gestalten wurde an dieser Stelle der Makro nicht mehr weiterentwickelt. 

\subsubsection{Endresultat}
\label{Endresultat}

Wenn der Makro erfolgreich durchlaufen ist, werden je nach Benutzereingabe eine oder zwei Matrizen ausgegeben. \\
Die erste Matrix wird immer erstellt und besteht aus dem Temperaturkollektiv der Messdateien. \\
die zweite Matrix wird nur erstellt, wenn die Information in der Messung enthalten und diese Information für die weitere Entwicklung benötigt wird. Da diese Matrix den Großteil der Berechnungsdauer in Anspruch nimmt, wäre es redundant diese Berechnung durchzuführen wenn sie nicht benötigt wird.\\
Da dieser Gesamtprozess wie oben aufgeführt ineffizient ist, bestand einer der weiteren Aufgaben darin diesen Prozess in MATLAB zu verbessern und den Workflow damit komplett zu überarbeiten.

Nach der Erfahrung des Authors nimmt beim Entwicklungsprozess die Datenaufarbeitung einen großen Anteil der gesamt benötigten Zeit und damit einen Großteil der Kosten auf sich. Die Kosten stellen sich aus Lizenzkosten der benötigten Programme und den direkten Personalkosten zusammen. \\
Daher kann mit einer entsprechenden Workflow-Verbesserung auch eine Kostensenkung erwirkt werden. \\

\newpage
\subsection{Aktueller Workflow der Messungsauswertung}

Der in \textbf{Datenauswertung in Excel} [\ref{Excel}] beschrieben Workflow beinhaltet nur einen Bruchteil der benötigten Arbeitschritte um verwertbare Ergebnisse aus dem Gesamtprozess zu gewinnen. \\

Da die Messdaten nur der letzte Teil einer Arbeitskette darstellen, wird nun auf den gesamtem Prozess eingegangen. Dieser wird erläutert, die Problemstellen hervorgehoben und eine Lösung in MATLAB für den Workflow vorgestellt. \\
%hier workflow diagramm einfügen

\subsubsection{Messstellenplan}

Der Messstellenplan oder \textbf{(MSP)} wird von den Messstellentechnikern und dem zuständigen Mitarbeiter erstellt und in Zusammenarbeit mit den relevanten Bauteilverantwortlichen und anderen Parteien kontinuierlich angepasst.\\
Die pro Versuchsreihe oftmals eine signifikante Anzahl von Messungen gefahren werden ist es recht aufwändig den Messstellenplan jedes Mal an die neue Messung anzupassen. \\
Der Messstellenplan wird über konzern-eigene Datenverwaltungstools für jeden Wagen gepflegt. Dieser kann dann in Excel exportiert werden und für jede relevante Messung referenziert werden. \\
Hierbei entstehen schnell eine unübersichtliche Anzahl von Excel Dateien mit zum Teil extensiven Bezeichnungen um eine genaue Zuordnung zu ermöglichen. \\

\subsubsection{Grenztemperaturen}

Die Grenztemperaturen werden in einer seperaten Datei für jede Messstelle, insofern diese Messstelle eine Grenztemperatur hat, festgelegt. \\
Analog zum \textbf{MSP} werden diese Grenztemperaturpläne in Excel exportiert. \\
Aktuell werden diese Dateien dann in einer seperaten Excel-Datei mit hinterlegten Makros eingelesen und ausgewertet. Da VBA eine recht alte und langsame Programmiersprache ist, ist diese Auswertung oft zeitaufwändig oder stürzt einfach ab. \\

\subsubsection{Messungen}

In der Excel-Datei können die Messungen aus ihrem nativen MDF-Format nicht eingelesen werden. \\
Daher müssen um die Messungen zuzuordnen erst die relevanten Kanäle mit Hilfe eines Porsche-internen Tools ausgelesen und in Excel exportiert werden. \\
Diese exportierten Excel-Dateie sind regelmäßig weiter über 100mb groß und benötigen dementsprechend lang um eingelesen zu werden. \\
Bei mehreren Messungen pro Auswertung ist hier der Aufwand schnell groß, da jede Messung einzeln verarbeitet werden muss. 

\subsubsection{Simulationen}

In dem aktuellen Workflow gibt es keine Möglichkeit Simulationsergebnisse der Auswertung anzufügen. \\
Da diese aber oft für Erstauslegungen oder Konzept-Bestätigung verwendet werden, wäre es von Vorteil diese direkt mit einsehen zu können.

\subsubsection{Ergebnisausgabe}

Die Ergebnisse werden entweder als PDF oder als Excel Workbook ausgegeben. \\
Diese werden dann oft in Powerpoint oder anderen Medien weiterverwendet um entsprechende Maßnahmen abzuleiten.

\newpage
\subsection{Workflow in MATLAB}

Um diesen Workflow effizienter zu gestalten wurde er in MATLAB übertragen und dort in einem Stand-alone Programm entworfen. \\

Der erste Arbeitschritt bestand darin, sich mit der MATLAB-App Designer Umgebung vertraut zu machen. Da MATLAB sehr effizient große Matrizen verarbeiten kann, bot sich das Programm für diese Art der Auswertung an. \\
Nach Durcharbeitung der Tutorials wurde als nächstes der Workflow entworfen. \\
Dieser sollte idealerweise die vorher genannte Auswertung in Excel und anderen Programmen komplett ersetzen. \\
%hier Workflow diagramm einfügen
Damit das Programm auch von den Kollegen eingesetzt wird, wurden die relevanten Parteien von Anfang an in den Entwicklungsprozess miteingebunden. Hierfür wurden wöchentliche Meetings angesetzt.\\
Im Laufe der Entwicklung des Workflows und Programms wurde schnell offensichtlich, wie wichtig diese andauernde Rückmeldung mit den anderen Teammitgliedern für eine erfolgreiche Implementierung ist.\\

\subsubsection{Vorarbeit}

Um den Workflow zu entwerfen musste erst der gesamte Auswertungsprozess verstanden und nachvollzogen werden. \\
Um dies zu erreichen wurden kleine Aufgaben entlang des gesamten Prozesses übernommen und selbstständig bearbeitet. Bei Unklarheiten wurden die relevanten Kollegen hinzugezogen.\\

Als nächstes wurde der Ablauf ausgelegt und dann Überlegungen angestellt, wo er beschleunigt und effizienter durhgeführt werden kann.

\subsubsection{Verbesserungen}

Die einfachste Verbesserung wurde im Bereich der Messstellenpläne implementiert. Diese können nun aus der exportierten Excel-Datei direkt in das Programm geladen werden. \\
Sollte eine neue Messung einen veränderten Messstellenplan haben, kann diese Veränderung in den Messstellenplan der bereits eingelesen ist eingefügt werden.\\
Um das weitere Vorgehen zu erleichtern und das Vergleichen von Messungen zu ermöglichen, werden neue Messstellen an den alten Plan angefügt und dann einsortiert. Keine Messstellen werden entfernt, damit alle vorhergehenden Messungen ihre Aussagekraft bewahren.\\

Des Weiteren wird der Grenztemperaturplan nur temporär eingelesen und die Werte für die relevanten Messstellen ausgelesen, wonach der Grenztemperaturplan wieder aus dem Arbeitsspeicher gelöscht wird. \\
Analog kann, sollte der Messstellenplan angepasst werden, die Grenztemperatur für die neuen Messstellen dann hinzugefügt werden. Sollte sich eine Grenztemperatur für eine bestehende Messstelle geändert haben, wird die alte Grenztemperatur automatisch überschrieben. Diese Maßnahme soll der Benutzerfreundlichkeit des Programms dienen. \\

Um die Messungen einzulesen, bedarf es nun keinem Dritt-Programm mehr. Das MATLAB-Programm gleicht den eingelesenen Messstellenplan in MDF-Format mit der gesamten Messung ab und lädt nur die Messkanäle die für den Plan von Relevanz sind.\\
Sollten weitere Kanäle wie zum Beispiel die GPS-Geschwindigkeit oder Ähnliches wichtig sein, können diese Messstellen manuell in den Messstellenplan in MATLAB eingefügt werden. Hierdurch werden diese Kanäle dann mitverwertet. \\
Genaueres hierzu in Kapitel \ref{Implementierung}.\\

Die Simulationsergebnisse liegen normalerweise in \textbf{.cvs} oder \textbf{.xml} Format vor. \\
Da MATLAB diese Formate nicht nativ einlesen kann, muss für diese Funktion noch eine Routine geschrieben oder eine verfügbare Funktion gefunden werden.\\
Aufgrund der wichtigeren Grundfunktion des Programmes wurde diese Funktion aber auf die 'Nice-To-Have'-Liste gesetzt. Daher wird an dieser Funktion nur gearbeitet, wenn der Rest des Programms zufriedenstellend läuft.\\

Das Ziel der Verbesserungen ist eine signifikante Zeiteinsparung und Vereinfachung des Workflows der Messungsauswertung und Ergebnisinterpretation.\\

\subsubsection{Implementierung und Probleme in der Entwicklung}
\label{Implementierung}
\subsubsection*{Version 0.1}

In der ersten Version des Programms wurde der Fokus darauf gesetzt die Grundfunktionen der Auswertung linear zu implementieren. \\
Dies sollte auch helfen den Grundprozess besser zu verstehen. Das Programm selbst sollte dann iterativ verbessert unde ausgebaut werden bis hin zur finalen Version.\\

Die ersten Probleme bei der Messstellenplanverwertung traten recht früh auf. Einige Messstellenpläne aus dem alten und dem neuen Datenverwaltungstool von Porsche waren zwar alles Excel-Dateien, aber mit verschiedenen Datentypen abgespeichert. \\
Zwischen diesen Datentypen gab es bei der Spaltennummerierung und -benennung unterschiede, was das auslesen der korrekten Information erschwert.\\
Nach Rücksprache mit dem Betreuer ergab sich, dass in kurzer Zeit das alte Datenverwaltungstool komplett durch das neue ersetzt werden würde. Daher wurde das Problem teilweise gelöst, jedoch bestand weiterhin in der Nomenklatur eine Diskrepanz. \\
Weiter Maßnahmen um die Spaltenbenennung zu vereinheitlichen sind daher nötig.\\
Meiner Meinung nach würde dies den Workflow verbessern, da man bei der Auswertung, ob manuell oder per Programm, immer mit der gleichen Informationsstruktur konfrontiert ist. \\
Die endgültige Lösung des Problems bestand dann darin, die Spalten mit den entsprechenden Namen einzulesen und in einer vordefinierten Reihenfolge zu speichern.\\
Als weitere gewünschte Funktion sollte der MSP direkt in MATLAB bearbeitet werden können. Da die grafische Umgebung von MATLAB in diesem Fall eher langsam arbeitet, wurde abgeraten dies im Programm zu machen. \\
Die Möglichkeit Messstellen umzubenennen, zu löschen oder hinzuzufügen wurden dennoch grafisch implementiert. 

Um den Grenztemperaturplan einzulesen, der auch in Excel exportiert wird, konnten die bereits implementierten Methoden wiederverwendet werden. \\
Analog zum Messstellenplan ist das hinzfügen, löschen und umbenennen oder editieren der Grenztemperaturen und Messstellen hier möglich. Dafür wird die eingelesene Information in einer Tabelle dargestellt. \\

Die signifikanteste Zeiteinsparung wird bei dem Auswerten der Messungen erzielt. \\
Um die Messungen einzulesen wurden Porsche-interne Routinen verwendet. Diese lesen aus den \textbf{.mf4}-Messdateien die relevanten Informationen wie Zeitskala, Werte, Messstellennamen und -info aus. \\
Diese Informationen werden in einem Cell-Array gespeichert und können so abgerufen werden. \\
Dann gleicht eine Funktion die Messstellennamen mit dem Messstellenplan und Grenztemperaturplan ab und speichert alle gefundenen Übereinstimmungen in ein weiteres Cell-Array, in dem auch der MSP mit den Grenztemperaturen ist, ab. \\
Dieses Array ist vordefiniert als 3-D Array, welches in Richtung der Messstellen, oder der ersten Dimension, dynamisch erweitert werden kann. \\
Obwohl das dynamische Erweitern von Arrays in MATLAB die Rechendauer verlängert, war dies die effizienteste Lösung.\\
Die zweite Dimension ist definiert unde beinhaltet alle relevanten Daten der einzelnen Messstellen, sowie die jeweiligen Messstellen. \\

\newpage
Die dritte Dimension beinhaltet die Messdaten der Messstellen. Diese Dimension wurde als \textit{k} vordefiniert mit 

\begin{equation}
	k = 100
\end{equation}

Da MATLAB sonst die zum Teil aus mehreren Hundert Messstellen bestehenden Messungen dynamisch hinzufügen müsste, was zu einer signifikanten Verlängerung der Einlesezeit geführt hätte.\\
Vor dem Einlesen jeder Messung wird der Benutzer gefragt, ob die Messung neue Messstellen beinhaltet. Diese können dann bei Bedarf aus Excel eingelesen und zu den existierenden Messstellen hinzugefügt werden. \\
Die Messinformation der Messungen werden in einem zweiten Array abgespeichert, dass mit den gleichen Indizes arbeitet. So können jederzeit alle Information zusammen mit den Messungen abgerufen werden. \\

\subsubsection{Auswertung}

Für die Auswertung müssen nach dem einlesen aller Daten noch einige Parameter vom Benutzer angepasst werden, je nachdem welche Informationen betrachtet werden sollen.\\
Dann wird aus den Daten eine PDF-Datei ertellt, die alle Messkurven mit den Informationen abbildet. \\

\subsubsection{Schwierigkeiten in MATLAB}

Das größte Problem mit MATLAB ist die langsame grafische Umgebung. \\
Da jedoch die Konsolenumgebung schwieriger zu bedienen ist, wurde die grafische Lösung gewählt.\\

Eine weitere Schwierigkeit stammt aus dem Erstellen einer Speicherungsdatei. Um nicht bei jeder Teilauswertung alle Dateien neu einlesen zu müssen, ist ein weiterer Schritt die Erstellung einer solchen Speicherdatei zu ermöglichen. \\
Hierfür müssen alle Daten die in der GUI hinterlegt sind ausgelesen und abgespeichert werden.\\
Dies wird bearbeitet, wenn am Ende des Projekts noch Zeit übrig ist. \\







		\cleardoublepage
	%\chapter{Schlussfolgerung und Ausblick}
\label{Conclusion}
%Was wurde in dieser Arbeit gemacht? Was sind die wesentlichen Ergebnisse? Wo ist weiterer Forschungsbedarf? Welche interessanten Forschungsbereiche ergeben sich aus der eigenen Arbeit?

In dieser Arbeit wurde eine Parameterstudie zur Untersuchung von Stößen auf isotropen Platten durchgeführt. Zuerst wurde die benötigte Theorie in Kapitel~\ref{chap:Principles} hergeleitet. Anschließend wurden im Rahmen der Studie die Parameter in Tabelle~\ref{tab:VariierteParams} ausgehend vom Ausgangsfall in Tabelle~\ref{tab:Ausgang} variiert. In Kapitel~\ref{chap:Implementation} wurden die Veränderungen im Skript dargestellt und die Funktionsweise des Programms erklärt. 

\begin{table}[H]
	\begin{center}
		\caption{Parameterstudie}
		\label{tab:VariierteParams}
		\begin{tabular}{l|c}
			\textbf{Variable} & \textbf{Wert}\\
			\hline
			$h$ & Höhe der Platte [cm]\\
			$v_{0}$ & Geschwindigkeit der Impaktors [cm/s]\\
			$r_{g}$ & Radius des Impaktors [cm]\\
			$\frac{\xi}{a}$ & Auftreffstelle in x-Richtung [entdimensioniert]\\
			$\frac{\eta}{b}$ & Auftreffstelle in y-Richtung [entdimensioniert]\\
			$Sv$ & $\frac{a}{b} \equiv \; \mbox{Seitenverhältnis der Platte}$ \\		
		\end{tabular}
	\end{center}
\end{table}

Wie in Kapitel~\ref{chap:Durchfuehrung} ausgeführt, lassen sich einige Schlussfolgerungen aus den gewonnenen Ergebnissen ziehen. \\

Im Wesentlichen lässt sich sagen, dass die Plattensteifigkeit einen großen Einfluss auf das Verhalten der Platte hat. Allgemein lässt sich feststellen, dass eine erhöhte Plattensteifigkeit zu einer geringen Anzahl von Aufschlägen führt.

Eine Vergrößerung der Höhe der Platte zu einer erhöhten Steifigkeit und somit zu einer geringeren Durchbiegung und erhöhten maximalen Kraft. Hier liegt ein Kraftunterschied von $P(Mr_{max}) \; \approx \; 8 \cdot P(Mr_{min})$ vor. Dies ist auf die Potenz der Höhe in $\overline{a}$ in Kapitel~\ref{chap:Principles} zurückzuführen.

Wenn das Seitenverhältnis verändert wird, dominiert die kürzere der beiden Seiten das Verhalten der Platten. Zunächst steigt ebenfalls die maximale Kraft, sowie die maximale Durchbiegung der Platte. Ab einem Seitenverhältnis von $2.5$ fallen jedoch beide kontinuierlich ab.

Betrachtet man die Geschwindigkeit des Impaktors, wird ersichtlich, dass hier $v_{0}$ der dominierende Parameter ist. Zwar steigt die Auslenkung mit größer werdendem $Mr$ an, jedoch ist diese Änderung wesentlich kleiner als der Anstieg bei konstantem $Mr$ und größer werdender Geschwindigkeit. Die maximale Auslenkung und maximale Kraft zeigen einen linearen Zusammenhang bei einer Veränderung der Geschwindigkeit des Impaktors. Somit führt eine erhöhte Geschwindigkeit stets zu einer erhöhten Kraft und einer erhöhten maximalen Durchbiegung.

Der Radius der Kugel hat nur einen verschwindend geringen Einfluss auf die Auslenkung der Platte. Wie in Kapitel~\ref{chap:Durchfuehrung}, Abbildung~\ref{fig:RadiusAuslenkung} zu sehen ist, steigert sich die Auslenkung von $r_{g,min}$ zu $r_{g,max}$ nur in der zweiten Nachkommastelle.
Betrachtet man die maximal auftretende Kraft, sieht man, dass hier die Höhe und die Geschwindigkeit den größten Einfluss haben. In Abbildung~\ref{fig:Faktoren} sind die Faktoren von $Mr_{min}$ und $Mr_{max}$ bei den Grenzwerten der Parameter dargestellt. In Abbildung~\ref{fig:Faktoren}(b) kann man ablesen, dass die Faktoren des Radius, analog zur Auslenkung, zwischen der minimalen und maximalen Kraft bei $Mr_{min}$ und $Mr_{max}$, verglichen mit den Faktoren der Höhe und der Geschwindigkeit, gering sind. \\

Einen großen Einfluss auf die maximale Kraft und Durchbiegung zeigt der Aufschlagsort.
Die Anzahl der Schläge ist direkt mit dem Abstand zum Rand der Platte beziehungsweise mit dem Abstand zum Mittelpunkt der Platte verknüpft. Mit steigender Steifigkeit zum Rand hin sinkt die maximale Anzahl an Schlägen. Jedoch sind weder die maximale Kraft noch die maximale Durchbiegung in der Mitte der Platte vorzufinden, diese befinden sich in einem ringförmigen Gebiet welches einen Radius von $0.2 \cdot a$ für den Fall einer quadratischen Platte hat. Auffallend ist, dass die geringste Kraft in der Plattenmitte vorzufinden ist. Auffällig ist jedoch, dass die Durchbiegung in der Plattenmitte, verglichen zu der Durchbiegung bei $\frac{\xi}{a} = \frac{\eta}{b}= 0.75$ geringer ausfällt.


\begin{figure}[H]%
	\centering
	\subfloat[Faktoren der Auslenkung]{{\includegraphics[width=0.4\linewidth]{pictures/gnuplot/3d/Hoehe/production/Auslenkungsfakt.eps} }}%
	\qquad
	\subfloat[Faktoren der Kraft]{{\includegraphics[width=0.4\linewidth]{pictures/gnuplot/3d/Hoehe/production/Kraftfakt.eps} }}%
	\caption{Faktoren: (v.L.n.R) Höhe, Geschwindigkeit und Radius}%
	\label{fig:Faktoren}%
\end{figure}

\newpage

\section*{Ausblick}
\label{sec:Ausblick}

Wie in dieser Arbeit beschrieben, haben vor allem die Plattenhöhe, die Auftreffgeschwindigkeit und die Auftreffstelle einen großen Einfluss auf das beschriebene Problem. Bei der Schadensbeurteilung sollten diese drei Faktoren berücksichtigt werden. Auch das Seitenverhältnis $Sv$ hat einen signifikanten Einfluss, jedoch sind die Ergebnisse durch die auftretenden numerischen Instabilitäten nur bedingt verlässlich. In diesem Bereich sollte in zukünftigen Arbeiten weitere Recherche betrieben werden. \\
Ein weiteres Forschungsthema liegt in dem Bereich der dünnen Platten mit $h \leq 0.5 cm$. Eine Verfeinerung der Schrittweite könnte Platten mit $h \leq 0.5 cm$ auflösen, jedoch würde der zeitliche Anspruch deutlich steigen. Da oftmals dünne Platten aus Gewichtsgründen verbaut werden, sind hier weitere Überlegungen gefragt. \\
Ebenfalls wurde weder auf die maximale Spannung noch auf Quer- und Membrankräfte in der Platte eingegangen, so dass hier weitere Recherchen notwendig sind.

Wie bereits in der Auswertung des Seitenverhältnisses ersichtlich, wurden Membrankräfte vernachlässigt. Dies führt zu deutlich höheren Durchbiegungen als die Realität bieten könnte.
Ein nicht-linearer Ansatz könnte dieses Problem beheben, jedoch würde die Berechnung beispielsweise durch gegebene Finite-Elemente Methoden erheblich mehr Zeit in Anspruch nehmen und qualitativ die gleichen Lösungen aufzeigen.

 	\cleardoublepage
%
% 			BIBLIOGRAPHY
%---------------------------------------------------------------------------------------
\addcontentsline{toc}{chapter}{Literatur}									% Literaturverzeichnis in Inhaltsvz aufnehmen
\nocite{*} 																	% Alle Literatureintraege ins Vz. aufnehmen
																			% BITTE AUSKOMMENTIEREN FUER DIE ARBEIT
\renewcommand{\bibname}{Literaturverzeichnis}								% NUR ZU ANSCHAUUNGSZWECKEN IN DEM BEISPIEL
\printbibliography				\cleardoublepage							% Literaturverzeichnis 
%
% 			APPENDIX
%---------------------------------------------------------------------------------------
\appendix
\pagestyle{scrheadings}														% Koma-script Kopf-/Fusszeile , Definition s.o.
	%\chapter{Anhang - Tipps zum Verfassen studentischer Arbeiten}
\minisec{Allgemeine Regeln:}
\begin{compactitem}
\item[\emph{Neutrale Formulierung:}] Die Arbeit wird in der 3. Person Singular verfasst.
\item[\emph{Tempus:}] Die Arbeit wird durchgängig im Präsens geschrieben.
\item[\emph{Aktive Formulierung:}] Formulierungen im Passiv vermeiden.
\item[\emph{Formulierung im Indikativ:}] Die Verwendung des Konjunktivs vermeiden.
\item[\emph{Abkürzungen:}] Alle Abkürzungen, die nicht im Duden stehen, müssen eingeführt werden.
\end{compactitem}
\minisec{Generell gilt:}
\begin{itemize}
\item[\emph{So wenig wie möglich, so viel wie nötig!}] Herleitungen und Erklärungen auf das Wesentliche beschränken und nicht zu weit abschweifen; Vermeidung der Aufzählung unnötiger Einzelheiten, von allgemein Bekanntem und von Informationen, die nicht zum Thema gehören.
\item[\emph{So einfach wie möglich:}] Möglichst kurz und prägnant formulieren. 
\item[\emph{Kurze Sätze:}] Vermeidung von komplizierten Satzreihen und Schachtelsätzen sowie von langen zusammengesetzten Wörtern.
\item[\emph{Fachvokabular Sätze:}] An den richtigen Stellen einheitlich und durchgängig verwenden. Für das Verständnis notwendige, aber nicht geläufige Fach- und Fremdwörter sowie Abkürzungen erläutern.
\item[\emph{Redundanz:}] Unbekannte und besonders wichtige Informationen müssen dem Leser durch angemessene Wiederholung vors Auge geführt werden. 
\item[\emph{Übersichtlichkeit:}] Ergebnisse, Versuchspläne, Maschinendaten, etc. möglichst in Tabellenform abfassen und nicht im Fließtext aufzählen.
\item[\emph{Aussagekräftige Kapitelüberschriften}] (nicht zu lang!) erleichtern das Verständnis.
\item[\emph{Unterkapitel}] sind nur bis zur vierten Stufe zulässig; Eine neue Stufe ist nur dann einzuführen, wenn es mindestens zwei Kapitel in dieser Ebene gibt.
\item[\emph{Kapiteleinleitung und -abschluss:}] Ein bis zwei \emph{kurze Einleitungssätze} am Anfang von jedem Kapitel (z.B. In diesem Kapitel geht es um..., Im Folgenden wird beschrieben...) bzw. \emph{zusammenfassende Sätze} am Ende eines Kapitels geben dem Text eine Struktur und tragen so zum allgemeinen Verständnis bei.
\end{itemize}

Der Umfang der schriftlichen Arbeit richtet sich nach der jeweiligen Prüfungsordnung. Es gilt:
\begin{center}
\begin{tabular}{lr}
\toprule
{Projektarbeit} & k.A.\\
{Bachelorarbeit} & < 50 Seiten ohne Anhang\\
{Masterarbeit} & < 80 Seiten ohne Anhang\\
\bottomrule
\end{tabular}
\end{center}
%
\minisec{Kontrollfragen:}
\begin{itemize}
\item Ist der \emph{Rote Faden} durch die ganze Arbeit erkennbar?
\item Folgt der Aufbau der Arbeit einer \emph{logischen Struktur} (Reihenfolge, Gedankensprünge, Sinnabschnitte,etc.)?
\item Gibt es eine durchgängige \emph{Zeitform} (Präsens), oder gibt es Zeitsprünge (z.B. Präteritum im Durchführungsteil)?
\item Stimmen die \emph{inhaltlichen Bezüge}? (Passt die Verbform zum Subjekt?)
\item Sind die Inhalte \emph{übersichtlich} gestaltet?
\end{itemize}


			\cleardoublepage
	%\chapter{Anhang - Formale Fragen}
Die Formatierung der Vorlage ist beizubehalten. Die Farben in Diagrammen und Zeichnungen müssen so gewählt werden, dass durch das Ausdrucken der Arbeit in Schwarz-Weiß keine Informationen verloren gehen (Graustufen). Die Verwendung der Farbe gelb ist unzulässig. Weiterhin muss auf Personen mit einer Rotgrünblindheit bei der Farbwahl Rücksicht genommen werden. Die Zahlen von eins bis einschließlich zwölf werden ausgeschrieben. Höhere Zahlen werden \emph{nicht} ausgeschrieben. Diagramme, Tabelle, etc. sind so zu formatieren, dass sie inklusive ihrer Unterschriften bei einer Verkleinerung auf DIN A5 lesbar bleiben.
%
\section{Umgang mit Latex - Beispiele}
%
\minisec{Formeln}
Die Theorie beinhaltet Formeln, Herleitungen, empirische Daten, etc. Für die Darstellung von Formeln in der Formelumgebung bzw. im Fließtext gelten einige Regeln:

Das Multiplikationszeichen ist der Punkt und nicht der Stern. Verwendet wird er bei Skalarprodukten, oder wenn sonst Irrtümer entstehen können. Ansonsten macht ein Multipliaktionszeichen Formeln unübersichtlich lang. Um die Eindeutigkeit und Übersichtlichkeit in Formeln zu gewährleisten bietet es sich häufig an zusätzliche Abstandshalter und verschieden große Klammern zu verwenden. 
\begin{equation}
\begin{split}
%\dot{\gls{x}}&~=~\gls{A}\gls{x}(t)~+~%
%							  \gls{B}\gls{u}(t)\\
%\gls{y}&~=~\gls{C}\gls{x}(t)~+~\gls{D}\gls{u}(t)\\
\dot{\underline{x}}&~=~\underline{A}\underline{x}(t)~+~%
							  \underline{B}\underline{u}(t)\\
\underline{y}&~=~\underline{C}\underline{x}(t)~+~\underline{D}\underline{u}(t)\:.
\end{split}
\label{eq:Gleichung1}
\end{equation} \glsadd{A}\glsadd{B}\glsadd{C}\glsadd{D}\glsadd{x}\glsadd{u}
\begin{equation}
\begin{pmatrix} 		
	\delta\dot{q_k}		\\
	\delta\dot{\alpha_k}	\\
	\delta\dot{V_k}		\\
	\delta\dot{\gamma_k} \end{pmatrix}%
=\underbrace{\begin{bmatrix}	
	M_q			& M_{\alpha} 	& M_u 	& 0	\\ 
	1+Z_q 		& Z_{\alpha} 	& Z_u	& 0	\\
	X_q			& X_{\alpha}-g	& X_u 	& -g\\
	-Z_q 		& -Z_{\alpha} 	& -Z_u	& 0	\end{bmatrix}}_{\text{Zustandsmatrix}}%
\begin{pmatrix} 		
	\delta{q_k}			\\
	\delta{\alpha_k}	\\
	\delta{V_k}			\\
	\delta{\gamma_k} 	\end{pmatrix}%
+\underbrace{\begin{bmatrix}
	M_{F} 		& M_{\kappa}	& M_{\eta}\\
	Z_{F} 		& Z_{\kappa}	& Z_{\eta}\\
	X_{F} 		& X_{\kappa}	& X_{\eta}\\
	-Z_{F} 		& -Z_{\kappa}	& -Z_{\eta}\end{bmatrix}}_{\text{Eingangsmatrix}}%
\begin{pmatrix}
	\delta_F		\\
	\delta\kappa	\\
	\delta\eta 		\end{pmatrix}\:.
\label{eq:gleichung2}
\end{equation}
In der Formel-Umgebung wird Text so dargestellt, wie im Fließtext. Dazu verwendet man den Befehl $\backslash$text oder $\backslash$mathrm. Letzterer empfiehlt sich u.a. für Einheiten im Formeltext, da lediglich die Schriftart geändert wird, die Formel-Umgebung jedoch erhalten bleibt. Auch der Differentialoperator (Nicht zu verwechseln mit der partiellen Ableitung) bzw. Integrale werden nicht kursiv geschrieben:
\begin{equation}
E=\frac{\operatorname{d} \! \sigma}{\operatorname{d} \! \epsilon} \:.
\end{equation}
Partielle Ableitung:
\begin{equation}
\tau=G \;\left(\frac{\partial u}{\partial y} + \frac{\partial v}{\partial x}\right) \:.
\end{equation}

In abgesetzten Formeln werden Brüche mit horizontalem Bruchstrich (Befehl $\backslash$frac{}{}) verwendet  $\mathrm{\frac{N}{mm^2}}$. Im Fließtext werden Brüche mit einem Schrägstrich dargestellt. Dabei ist ausdrücklich die Formatierung \emph{ohne} Fraction-Umgebung, $\mathrm{N/mm^2}$, der Formatierung \emph{mit} einer Umgebung, wie dem Befehl $\backslash$sfrac{}{}, $\mathrm{\sfrac{N}{mm^2}}$, vorzuziehen.

Die Zeichen für die Dezimal- und Tausendertrennung richten sich nach der Sprache der Arbeit (vgl. Tabelle \ref{tab:Trennzeichen}). In Grafiken sollte man sich an die jeweiligen Konventionen halten, in deutschsprachigen Arbeiten ist es jedoch ebenso zulässig, die englische Konvention zu nutzen, sofern eine andere Darstellung nicht mit vertretbarem Aufwand erreichbar ist.
\begin{table}
\centering
\begin{tabular}{lrr}
\toprule
							& Deutsch	& Englisch\\
\midrule
Tausendertrennzeichen		& .			& ,\\
Dezimaltrennzeichen 		& ,			& .\\
\bottomrule
\end{tabular}
\caption{Dezimal- und Tausendertrennung, aufgeführt nach Sprachraum.}
\label{tab:Trennzeichen}
\end{table}
%
\minisec{Abbildungen}
Das Hinzufügen von Grafiken und Abbildungen dient der Veranschaulichung bzw. der Ergänzung des Fließtexts. Abbildungen werden mit der $\backslash$figure-Umgebung in den Text eingebunden. Der Befehl mit dem die jeweilige Abbildung aufgerufen wird lautet:\\
$\backslash$includegraphics[ scale=.71]\{pictures/Dateiname.eps\}. Dazu muss die verwendete Datei als .eps- bzw. als .png-Datei im Ordner \glqq pictures\grqq \! abgelegt werden. Mit Befehlen wie beispielsweise scale, angle, width, oder height in einer Latex-spezifischen Einheit wird das Bild bezüglich Skalierung, Drehwinkel, Breite oder Höhe angepasst. Bei Zeichnungen ist darauf zu achten, dass Zeichnungs- und Bemaßungslinien, sowie Schriftbreiten mit einem normgerechten Satz von Strichbreiten erstellt wurden. Bei Verwendung des Strichbreiten-Satzes [0,35mm; 0,7mm] ist in der Zeichnung eine Schriftgröße von 16pt in etwa normgerecht. Eine Skalierung dieser Zeichnung mit dem Faktor $\mathrm{=\:1/\sqrt{2}}$ ergibt einen Strichbreitensatz von [0,25mm; 0,5mm] und die entsprechende Schriftgröße von 12pt. Ein Beispiel ist in Abbildung \ref{fig:Bsp_Zeichung} dargestellt.
\begin{figure}[t!]
\centering
\includegraphics[scale=0.71]{pictures/Beispiel.eps}
\caption{Dieses Beispielbild wurde mit dem normgerechten Satz der Strichbreiten [0,35mm; 0,7mm] und einer Schriftgröße von 16pt erstellt und gemäß der gewählten Schriftgröße im Fließtext von 12pt mit dem erforderlichen Faktor von scale=.71 ($\mathrm{=\:1/\sqrt{2}}$) skaliert.}
\label{fig:Bsp_Zeichung}
\end{figure}
%
\minisec{Tabellen}
Tabelle \ref{tab:refCase} zeigt ein Beispiel für eine Standard-Tabelle mit der Umgebung \emph{$\backslash$tabular}. Hier wird auch die Nutzung einiger weiterer typischer Befehle in Tabellen deutlich. Sollte eine Tabelle von sich aus nicht über die gesamte Breite gehen, bietet es sich aus optischen Gründen häufig an, diese zu über eine definierte Textbreite zu strecken. Dies erfolgt mithilfe des Pakets TabularX. Ein Beispiel ist in Tabelle \ref{tab:refTabX} gezeigt. Sollte jedoch eine Tabelle so schmal ausfallen, dass die Streckung unvorteilhaft wird, so kann eine einfache Tabular-Umgebung, wie in Tabelle \ref{tab:refCase_simple} gewählt werden.

\begin{table}[tb]
\centering
\begin{tabularx}{\textwidth-4em}{m{6mm}Xrrrrrrrr}
\toprule
&		& \multicolumn{8}{c}{Mach number}\\ 
\cmidrule{3-10}
& 		& 0.2	& 0.4	& 0.6	& 0.75	& 0.82	& 0.85	& 0.88& 0.96\\
\midrule %\cmidrule{3-10}
\multirow{7}{*}{\begin{sideways} dyn. pressure $\mathrm{[10^3 kg\:/m\:s^2]}$ \end{sideways}}%
& 2.00 	& 2800	& 12000	& 0		& 0		& 0		& 0		& 0		& 0	\\
& 2.30 	& 1700	& 11000	& 0		& 0		& 0		& 0		& 0		& 0	\\
& 2.70 	& 400	& 10000	& 0		& 0		& 0		& 0		& 0		& 0	\\
& 3.30 	&-1000	& 9300	& 14000	& 0		& 0		& 0		& 0		& 0	\\
& 4.10 	& 0		& 7600	& 13000	& 1500	& 0		& 0		& 0		& 0	\\
& 5.80 	& 0		& 5300	& 10000	& 1200	& 14000	& 15000	& 15000 & 0	\\
& 8.00 	& 0		& 2800	& 8700	& 1100	& 12000	& 13000	& 13000 & 14000\\
\midrule
\multicolumn{10}{c}{altitude [$\mathrm{m}$]}\\
\bottomrule
\end{tabularx}
\caption{Dies ist eine Tabellenunterschrift, die so lang ist, dass sie sich über zwei Zeilen erstreckt. Man erkennt, dass die nachfolgenden Zeilen eingerückt werden.}
\label{tab:refCase}
\end{table}
\begin{table}[bt]
\centering
\begin{tabularx}{\textwidth-4em}{Xrrr}
\toprule
				&$\omega_0~\mathrm{[rad/s]}$ &$\zeta~\mathrm{[-]}$& Polstellen\\
\midrule
Erste Biegemode - ungeregelt	& 10.0000	& 0.0700	& -0.6700 $\pm$ 10.0036i\\
Erste Biegemode - geregelt		& 9.9000	& 0.0750	& -0.7500 $\pm$ 9.9652i\\[1ex]
AS-Mode - ungeregelt			& 1.9000	& 0.5700	& -1.1000 $\pm$ 1.5600i\\
AS-Mode - geregelt 				& 1.9000	& 0.5700	& -1.1000 $\pm$ 1.5600i \\[1ex]
Phygoiden-Mode - ungeregelt		&			& 			& -0.2000 $\wedge$ 0.3000\\
Phygoiden-Mode - geregelt		&			& 			& -0.2000 $\wedge$ 0.3000\\
\bottomrule
\end{tabularx}
\caption{Dies ist eine Tabelle mit TabularX, deren Breite auf die Textbreite abzüglich je links und rechts 2em ausgestreckt wurde.}
\label{tab:refTabX}
\end{table}
Blindtext: Acto re stupeo Labor sus, ver ex aut exhorto sis aliter foetidus expono. Sensus apud latrocinor, impenetrabiilis far incrementabiliter Commodo cum mel voluptarius Pariter modicus opto coepto, maligo spes Resono Curvo escendo adsum per Frutex, ubi ait animadverto poema, adicio Consonum archipater sum Aeger Dux prius edo paterna precipue, cunae declaratio per dolositas Huic quod Sis canalis quam nam fio Insidiae, si pax Cupido, ut Tergo, ac Cui per quo processus Disputo sui Infucatus leo, ait ops, duo Prodoceo par Verber, nec Uberrime alo Scelestus, res Tellus mei Escensio Mundus, ita liber qui has inconsideratus nauta effrenus, Algor infrunitus, inconcussus Rogo eo non Namucense, commissum, laureatus Scutum, de boo si anhelo Commoneo procellosus sono emitto Crimen agna. Si subo Accubo castimonia hic ibi qua lux sto eu Pulcher Sem. Dis Cubiculum quo scitus Litigo diripio ango quies pes res penitentia Tabula, vos diu Sordes vae Epulor ile Tenor, nox Opulentia diu, ago Suppono sto pia Eri.
%
\minisec{Nützliche Befehle}
Mit dem Befehl $\backslash$emph werden wesentliche Begriffe \emph{hervorgehoben}. Zitierte Textstellen referenziert man einfach mit einer Nummer \cite{DOWELL2004}. Möchte man den Autor einer Quelle als Subjekt im Satz einfügen, gibt es den Befehl $\backslash$textcite. Ein Beispiel: \textcite{DOWELL2004} beschreiben wesentliche Aspekte. 

Absätze werden, wie bekannt aus MS Word, mit einer Leerzeile erzeugt. Einen Zeilensprung erhält man mit dem Befehl $\backslash\backslash$. Zur Veranschaulichung folgt ein Beispiel. Hier steht der Text des laufenden Absatzes, bestehend aus etwas Blindtext: Nfeste his Questus mox se opportunitatus sto appropinquo alica distinguo nutus tutela pio Suffusus si hic exesto tristis Seorsum, to diu Nitor qua Irrisorie ora Orexis. Dieser Satz beendet den laufenden Absatz mit dem Befehl $\backslash\backslash$.\\
Dieser Satz folgt auf den einfachen Zeilensprung mit dem Befehl $\backslash\backslash$. Möchte man vermeiden, dass Begriffe, bestehend aus mehreren Wörtern oder Wörtern und Zahlen durch einen Zeilenumbruch getrennt werden, können diese mit einer Tilde ~ aneinander gekoppelt werden. Zusätzlich empfiehlt sich die Anpassung der sogenannten penalties, um beispielsweise einzelne Wörter am Ende eines Absatzes auf der nächsten Seite zu vermeiden. Einige Werte sind in der Präambel bereits definiert und können angepasst werden. Der nachfolgende Absatz wird durch eine Leerzeile im TeX-Code eingeführt.

Um bestimmte Punkte kurz aufzulisten empfiehlt sich die Umgebung \emph{compactitem}. Im Vergleich zu der Umgebung \emph{itemize} rückt dieser Befehl die einzelnen Auzählungspunkte näher zusammen:
\begin{compactitem}
\item Punkt 1 beinhaltet etwas Text,
\item Punkt 2 ist auch interessant und
\item Punkt 3 ist ebenfalls ein wesentlicher Aspekt.
\end{compactitem}

Vor der Einleitung wird das Symbol- und Abkürzungsverzeichnis erstellt. Um Abkürzungen oder Symbole im Symbolverzeichnis zu referenzieren verwendet man die Befehle $\backslash$gls, bzw. $\backslash$glsadd, falls im Verzeichnis nur ein Verweis auf diese Seite erfolgen soll. Der Befehl $\backslash$gls erzeugt unter Verwendung des hyperref packages (mit der Option: draft = false) einen markierten Link in das Symbolverzeichnis. Im Symbolverzeichnis beschriebene Begriffe erscheinen erst, wenn diese im Text mittels einer der Befehle referenziert werden. Dies kann vollständig mit dem Befehl $\backslash$glsaddall[types=\{nomenclature, subscripts\}] umgangen werden. \glsadd{ROMAN}\glsadd{GREEK} \glsadd{eta}\glsadd{Theta}\glsadd{gamma}\glsadd{kappa}\glsadd{alpha} \glsaddall[types={subscripts,notation}] %

Kompiliert wird der Text mit der Befehlsfolge pdflatex biber.exe pdflatex.exe. Um das Symbolverzeichnis mit dem Package \emph{glossaries} lauffähig zu machen, bitte in die Erläuterung der Präambel schauen.

\begin{table}[bt]
\centering
\begin{tabular}{lrr}
\toprule
Werkstoff	& Dichte				& E-Modul\\
			& $\mathrm{[g/cm^3]}$	& $\mathrm{[N/mm^2]}$\\
\midrule
Stahl		& 7,85					& 210000\\
AlCuMg1 	& 2,80					& 71500\\
CFK			& 1,50					& 150000\\
\bottomrule
\end{tabular}
\caption{Tabellenunterschrift}
\label{tab:refCase_simple}
\end{table}				\cleardoublepage
	%\chapter{Anhang mit wesentlichem Code}
\label{chap:Anhang_Skript}
\lstinputlisting{skript/Karas.py}

			\cleardoublepage

				
%---------------------------------------------------------------------------------------
%========  E N D = D O C U M E N T  ===================================================%
%---------------------------------------------------------------------------------------
\end{document}
%
%
%
%% end of file
